pdfoutput=1
% Options for packages loaded elsewhere
\PassOptionsToPackage{unicode}{hyperref}
\PassOptionsToPackage{hyphens}{url}
\pdfoutput=1
\documentclass[
  12pt,
]{article}
\usepackage{xcolor}
\usepackage{amsmath,amssymb}
\setcounter{secnumdepth}{-\maxdimen} % remove section numbering
\usepackage{iftex}
\ifPDFTeX
  \usepackage[T1]{fontenc}
  \usepackage{textcomp} % provide euro and other symbols
\else % if luatex or xetex
  \usepackage{unicode-math} % this also loads fontspec
  \defaultfontfeatures{Scale=MatchLowercase}
  \defaultfontfeatures[\rmfamily]{Ligatures=TeX,Scale=1}
\fi
\usepackage{lmodern}
\ifPDFTeX\else
  % xetex/luatex font selection
\fi
% Use upquote if available, for straight quotes in verbatim environments
\IfFileExists{upquote.sty}{\usepackage{upquote}}{}
\IfFileExists{microtype.sty}{% use microtype if available
  \usepackage[]{microtype}
  \UseMicrotypeSet[protrusion]{basicmath} % disable protrusion for tt fonts
}{}
\makeatletter
\@ifundefined{KOMAClassName}{% if non-KOMA class
  \IfFileExists{parskip.sty}{%
    \usepackage{parskip}
  }{% else
    \setlength{\parindent}{0pt}
    \setlength{\parskip}{6pt plus 2pt minus 1pt}}
}{% if KOMA class
  \KOMAoptions{parskip=half}}
\makeatother
\usepackage{longtable,booktabs,array}
\usepackage{calc} % for calculating minipage widths
% Correct order of tables after \paragraph or \subparagraph
\usepackage{etoolbox}
\makeatletter
\patchcmd\longtable{\par}{\if@noskipsec\mbox{}\fi\par}{}{}
\makeatother
% Allow footnotes in longtable head/foot
\IfFileExists{footnotehyper.sty}{\usepackage{footnotehyper}}{\usepackage{footnote}}
\makesavenoteenv{longtable}
\setlength{\emergencystretch}{3em} % prevent overfull lines
\providecommand{\tightlist}{%
  \setlength{\itemsep}{0pt}\setlength{\parskip}{0pt}}
\usepackage{bookmark}
\IfFileExists{xurl.sty}{\usepackage{xurl}}{} % add URL line breaks if available
\urlstyle{same}
\hypersetup{
  hidelinks,
  pdfcreator={LaTeX via pandoc}}

\author{SCX}
\date{2026}

\begin{document}

\section{SCX Self-Evolution: Convergence Rate Analysis (Spring-1 Rate
Tightening)}\label{scx-self-evolution-convergence-rate-analysis-spring-1-rate-tightening}

\begin{quote}
\textbf{Version}: 2026-06-28 \textbar{} \textbf{Status}: Theoretical
derivation (partial) \textbar{} \textbf{Prerequisite}: Documents 02, 05,
06, 10 \textbf{Purpose}: Derive expected convergence rates for the SCX
self-evolution system, characterizing the dependence on Lipschitz
constants, learning rate schedules, and feature dimension. Move beyond
almost-sure convergence to quantitative rates.
\end{quote}

\begin{center}\rule{0.5\linewidth}{0.5pt}\end{center}

\subsection{Table of Contents}\label{table-of-contents}

\begin{enumerate}
\def\labelenumi{\arabic{enumi}.}
\tightlist
\item
  \hyperref[1-setup-and-notation]{Setup and Notation}
\item
  \hyperref[2-current-state-almost-sure-convergence-only]{Current State:
  Almost-Sure Convergence Only}
\item
  \hyperref[3-rate-under-strong-convexity-ot-alpha-derivation]{Rate
  Under Strong Convexity: \(O(t^{-\alpha})\) Derivation}
\item
  \hyperref[4-dependence-on-system-parameters]{Dependence on System
  Parameters}
\item
  \hyperref[5-rate-under-polyak-lojasiewicz-condition]{Rate Under
  Polyak-Łojasiewicz Condition}
\item
  \hyperref[6-rate-under-general-non-convexity]{Rate Under General
  Non-Convexity}
\item
  \hyperref[7-finite-time-non-asymptotic-bound]{Finite-Time
  (Non-Asymptotic) Bound}
\item
  \hyperref[8-two-timescale-rate-decomposition]{Two-Timescale Rate
  Decomposition}
\item
  \hyperref[9-comparison-with-known-rates]{Comparison with Known Rates}
\item
  \hyperref[10-gap-between-upper-and-lower-bounds]{Gap Between Upper and
  Lower Bounds}
\item
  \hyperref[11-summary]{Summary}
\end{enumerate}

\begin{center}\rule{0.5\linewidth}{0.5pt}\end{center}

\subsection{1. Setup and Notation}\label{setup-and-notation}

\subsubsection{1.1 The Coupled System}\label{the-coupled-system}

We analyze the convergence rate of \((S_t, \theta_t)\) to a fixed point
\((S^*, \theta^*)\) under the SCX self-evolution dynamics (Document 06):

\[\begin{aligned}
\theta_{t+1} &= \theta_t - \alpha_t \nabla_\theta \ell(f_{\theta_t}(x_t), y_t), \quad (x_t, y_t) \sim P_{S_t}, \\
S_{t+1} &= \Pi_{[0,1]}[S_t + \beta_t (\text{SCXUpdate}(S_t, M_{t+1}, f_{\theta_{t+1}}) - S_t)].
\end{aligned}\]

\subsubsection{1.2 Key Parameters}\label{key-parameters}

\begin{longtable}[]{@{}
  >{\raggedright\arraybackslash}p{(\linewidth - 4\tabcolsep) * \real{0.3143}}
  >{\raggedright\arraybackslash}p{(\linewidth - 4\tabcolsep) * \real{0.2571}}
  >{\raggedright\arraybackslash}p{(\linewidth - 4\tabcolsep) * \real{0.4286}}@{}}
\toprule\noalign{}
\begin{minipage}[b]{\linewidth}\raggedright
Parameter
\end{minipage} & \begin{minipage}[b]{\linewidth}\raggedright
Meaning
\end{minipage} & \begin{minipage}[b]{\linewidth}\raggedright
Typical Range
\end{minipage} \\
\midrule\noalign{}
\endhead
\bottomrule\noalign{}
\endlastfoot
\(L_f\) & Lipschitz constant of \(f_\theta\) in \(\theta\) &
\(O(1)\)--\(O(10^2)\) \\
\(L_g\) & Lipschitz constant of \(\nabla \ell\) &
\(O(L_f \cdot L_\ell)\) \\
\(L_S\) & Lipschitz constant of SCXUpdate & \(O(1)\) \\
\(d_\theta\) & Student parameter dimension & \(10^4\)--\(10^8\) \\
\(d_\phi\) & Feature dimension & \(10^1\)--\(10^3\) \\
\(\mu\) & Strong convexity parameter (if applicable) &
\(0\)--\(O(1)\) \\
\(\alpha_t\) & Student learning rate & \(t^{-a}\), \(a \in (0.5, 1]\) \\
\(\beta_t\) & Gatekeeper update rate & \(t^{-b}\), \(b \in (a, 1]\) \\
\(\sigma_\xi^2\) & Gradient noise variance & Problem-dependent \\
\(G\) & Gradient bound & \(O(L_f \cdot B)\) \\
\(B_S\) & Gatekeeper update bound & \(1\) (scores in \([0,1]\)) \\
\end{longtable}

\subsubsection{1.3 Target Quantity}\label{target-quantity}

We seek bounds on:

\[\mathbb{E}\bigl[\|S_t - S^*\|_{M_0}^2\bigr] \quad \text{and} \quad \mathbb{E}\bigl[\|\theta_t - \theta^*\|^2\bigr],\]

where \(\|S\|_{M_0}^2 = \frac{1}{|M_0|}\sum_{x \in M_0} S(x, y(x))^2\)
is the empirical \(L^2\) norm on the reference set, and
\((S^*, \theta^*)\) is a fixed point.

\begin{center}\rule{0.5\linewidth}{0.5pt}\end{center}

\subsection{2. Current State: Almost-Sure Convergence
Only}\label{current-state-almost-sure-convergence-only}

\subsubsection{2.1 What Document 06
Establishes}\label{what-document-06-establishes}

Theorem SE-1 (Document 06) asserts \textbf{almost-sure convergence} to a
fixed point under conditions C1-C7, using a Lyapunov supermartingale
argument. The proof (Lemma SE-1.1) establishes:

\[\sum_{t=1}^\infty \alpha_t \|\nabla L_t(\theta_t)\|^2 < \infty \quad \text{a.s.}, \qquad \sum_{t=1}^\infty \beta_t \|\Delta S_t\|^2 < \infty \quad \text{a.s.}\]

These imply convergence to a stationary point, but provide \textbf{no
rate} --- they only guarantee that the gradient norms vanish in the
Cesàro sense, not how fast.

\subsubsection{2.2 What Rates Are Missing}\label{what-rates-are-missing}

\begin{longtable}[]{@{}
  >{\raggedright\arraybackslash}p{(\linewidth - 2\tabcolsep) * \real{0.3750}}
  >{\raggedright\arraybackslash}p{(\linewidth - 2\tabcolsep) * \real{0.6250}}@{}}
\toprule\noalign{}
\begin{minipage}[b]{\linewidth}\raggedright
Desired
\end{minipage} & \begin{minipage}[b]{\linewidth}\raggedright
Current Status
\end{minipage} \\
\midrule\noalign{}
\endhead
\bottomrule\noalign{}
\endlastfoot
\(\mathbb{E}[\|\theta_t - \theta^*\|^2] \leq C t^{-\alpha}\) &
\textbf{Not proven} \\
\(\mathbb{E}[\|S_t - S^*\|_{M_0}^2] \leq C t^{-\alpha}\) & \textbf{Not
proven} \\
Finite-time (non-asymptotic) bound with probability \(1-\delta\) &
\textbf{Not proven} \\
Dependence of \(\alpha\) on \(L_f\), \(\alpha_t\), \(d_\phi\) &
\textbf{Not characterized} \\
\end{longtable}

\subsubsection{2.3 Why Almost-Sure is
Insufficient}\label{why-almost-sure-is-insufficient}

Almost-sure convergence does not tell us: 1. How many iterations are
needed to reach a given accuracy. 2. How the rate degrades as dimension
\(d_\phi\) grows. 3. Whether the convergence is fast enough to be
practically useful. 4. How to tune \(\alpha_t, \beta_t\) for optimal
convergence speed.

\begin{center}\rule{0.5\linewidth}{0.5pt}\end{center}

\subsection{\texorpdfstring{3. Rate Under Strong Convexity:
\(O(t^{-\alpha})\)
Derivation}{3. Rate Under Strong Convexity: O(t\^{}\{-\textbackslash alpha\}) Derivation}}\label{rate-under-strong-convexity-ot-alpha-derivation}

\subsubsection{3.1 Strong Convexity
Assumption}\label{strong-convexity-assumption}

\textbf{Assumption SC (Strong Convexity of Limiting Loss).} The limiting
expected student loss
\(\bar{L}(\theta) = \mathbb{E}_{(x,y) \sim P_{S^*}}[\ell(f_\theta(x), y)]\)
is \(\mu\)-strongly convex in a neighborhood of \(\theta^*\):

\[\bar{L}(\theta') \geq \bar{L}(\theta) + \langle \nabla \bar{L}(\theta), \theta' - \theta \rangle + \frac{\mu}{2} \|\theta' - \theta\|^2, \quad \forall \theta, \theta' \in \mathcal{B}(\theta^*, R).\]

\textbf{Justification.} Strong convexity holds for linear models with
\(\ell_2\) regularization. For neural networks, it holds only in a local
neighborhood of a strict local minimum (which is generic for
overparameterized networks). This is a \textbf{local} assumption, valid
after the system enters the basin of attraction.

\subsubsection{3.2 Student Rate Under Stationary
Distribution}\label{student-rate-under-stationary-distribution}

\textbf{Theorem 11.1 (Student Rate, Stationary Distribution ---
PROVEN).} Assume SC (strong convexity), C2 (Lipschitz student), C4 (RM
rates with \(\alpha_t = \alpha_0 t^{-a}\), \(a \in (0.5, 1)\)), and that
the training distribution is fixed at \(P_{S^*}\). Then:

\[\boxed{\;\mathbb{E}[\|\theta_t - \theta^*\|^2] \leq C_1 \cdot t^{-a} + C_2 \cdot t^{-2a}\;},\]

where:

\[C_1 = \frac{2\sigma_\xi^2 \alpha_0}{\mu}, \qquad C_2 = \frac{4 G^2 \alpha_0^2}{\mu^2}.\]

For \(a \in (0.5, 1)\), the dominant term is \(C_1 t^{-a} = O(t^{-a})\).

\emph{Proof.} Standard SGD analysis under strong convexity (Bach \&
Moulines, 2011; Needell et al., 2014). Let
\(V_t = \mathbb{E}[\|\theta_t - \theta^*\|^2]\). From the SGD update:

\[\begin{aligned}
V_{t+1} &= V_t - 2\alpha_t \mathbb{E}[\langle \theta_t - \theta^*, g_t(\theta_t) \rangle] + \alpha_t^2 \mathbb{E}[\|g_t(\theta_t)\|^2] \\
&\leq V_t - 2\alpha_t \mathbb{E}[\langle \theta_t - \theta^*, \nabla \bar{L}(\theta_t) \rangle] + \alpha_t^2 G^2.
\end{aligned}\]

By strong convexity:
\(\langle \theta_t - \theta^*, \nabla \bar{L}(\theta_t) \rangle \geq \mu \|\theta_t - \theta^*\|^2\).
Thus:

\[V_{t+1} \leq (1 - 2\mu \alpha_t) V_t + \alpha_t^2 G^2.\]

With \(\alpha_t = \alpha_0 t^{-a}\), this recurrence solves to
\(V_t = O(t^{-a})\). The precise constant follows from the theory of
stochastic approximation with polynomially decaying step sizes (Polyak
\& Juditsky, 1992). \(\square\)

\textbf{Status: PROVEN.} This is a standard result for SGD on strongly
convex objectives.

\subsubsection{3.3 Gatekeeper Rate Under Fixed
Memory}\label{gatekeeper-rate-under-fixed-memory}

\textbf{Theorem 11.2 (Gatekeeper Rate, Fixed Student --- PROVEN).}
Assume C3 (Lipschitz gatekeeper), the student is fixed at \(\theta^*\),
and the gatekeeper update uses \(\beta_t = \beta_0 t^{-b}\) with
\(b \in (0.5, 1]\). Then:

\[\boxed{\;\mathbb{E}[\|S_t - S^*\|_{M_0}^2] \leq C_S \cdot t^{-b}\;},\]

where:

\[C_S = \frac{2 B_S^2 \beta_0}{1 - L_S^2} \quad \text{(provided } L_S < 1 \text{ for contraction)}.\]

\emph{Proof.} The SCX update toward consensus is approximately a
contraction when the memory bank is informative. Formally, if
\(\text{SCXUpdate}(S) = \hat{C} + \varepsilon(S)\) where
\(\|\varepsilon(S)\| \leq L_S \|S - \hat{C}\|\) with \(L_S < 1\), then:

\[\|S_{t+1} - \hat{C}\|_{M_0} \leq (1 - \beta_t(1 - L_S)) \|S_t - \hat{C}\|_{M_0} + \beta_t \|\text{noise}\|.\]

This recurrence gives \(O(t^{-b})\) convergence. \(\square\)

\textbf{Status: PROVEN under contraction.} The condition \(L_S < 1\)
(SCXUpdate is a contraction toward consensus) is \textbf{not guaranteed}
in general --- it depends on the quality of the consensus signal. When
the consensus is noisy (small \(\Delta_s\) in Theorem 1), \(L_S\) may
exceed 1.

\subsubsection{3.4 Coupled Rate Under Two-Timescale
Separation}\label{coupled-rate-under-two-timescale-separation}

\textbf{Theorem 11.3 (Coupled Rate Under Strong Convexity ---
CONJECTURED).} Assume SC (strong convexity), C2-C7, C6' (two-timescale,
\(\beta_t = o(\alpha_t)\)), and that the fixed point \((S^*, \theta^*)\)
is locally stable. Then:

\[\boxed{\;\mathbb{E}\bigl[\|\theta_t - \theta^*\|^2 + \|S_t - S^*\|_{M_0}^2\bigr] \leq C_\theta \cdot t^{-a} + C_S \cdot t^{-b} = O(t^{-a})\;},\]

since \(b > a\) implies \(t^{-b} = o(t^{-a})\). The \textbf{student rate
dominates} the convergence speed.

\emph{Proof sketch (gaps acknowledged).} Under two-timescale separation
(Document 05, Section 7), the student converges on the fast timescale to
the quasi-stationary point \(\theta^*(S_t)\), and the gatekeeper
converges on the slow timescale to \(S^*\). The coupled rate is the
slower of the two rates, which is \(t^{-a}\) (the student rate) since
\(b > a\). The key gap is the coupling error: the student's optimization
target shifts by \(O(\beta_t)\) per step, introducing an additional
error term of order \(\beta_t / \alpha_t = t^{-(b-a)} = o(1)\). Under
strong convexity, this tracking error is bounded and does not affect the
asymptotic rate.

\textbf{Status: CONJECTURED.} The proof requires controlling the
accumulation of tracking errors \(\|\theta_t - \theta^*(S_t)\|\), which
is standard in two-timescale SA (Borkar, 2008, Chapter 6) but requires
verification for the specific SCX coupling structure.

\begin{center}\rule{0.5\linewidth}{0.5pt}\end{center}

\subsection{4. Dependence on System
Parameters}\label{dependence-on-system-parameters}

\subsubsection{\texorpdfstring{4.1 Dependence on Lipschitz Constant
\(L_f\)}{4.1 Dependence on Lipschitz Constant L\_f}}\label{dependence-on-lipschitz-constant-l_f}

The constant \(C_1\) in Theorem 11.1 depends on \(L_f\) through:

\[C_1 = \frac{2\sigma_\xi^2 \alpha_0}{\mu}, \quad \sigma_\xi^2 \leq G^2 = O(L_f^2), \quad \mu = \Omega(1/L_g) = \Omega(1/L_f).\]

\textbf{Result}: \(C_1 = O(L_f^3)\). A larger Lipschitz constant
(rougher loss landscape) degrades the convergence rate cubically: twice
the Lipschitz constant means 8× slower convergence.

\textbf{Corollary 11.1 (Lipschitz Dependence).} Under strong convexity
with \(L_g = \Theta(L_f)\):

\[\mathbb{E}[\|\theta_t - \theta^*\|^2] \leq \tilde{O}\!\left(\frac{L_f^3 \cdot \sigma_0^2}{\mu_0} \cdot t^{-a}\right),\]

where \(\sigma_0^2\) is the base noise level and \(\mu_0\) is the base
curvature.

\subsubsection{4.2 Dependence on Learning Rate
Schedule}\label{dependence-on-learning-rate-schedule}

The exponent \(\alpha\) in \(O(t^{-\alpha})\) equals the learning rate
exponent \(a\) (from \(\alpha_t = \alpha_0 t^{-a}\)).

\textbf{Proposition 11.1 (Optimal Rate Exponent).} Under strong
convexity, the optimal choice is:

\[a = 1 \quad \text{(harmonic learning rate } \alpha_t = \alpha_0 / t \text{)},\]

giving \(\mathbb{E}[\|\theta_t - \theta^*\|^2] = O(t^{-1})\).

However, this requires precise tuning of \(\alpha_0\): too small gives
slow convergence, too large causes divergence. In practice,
\(a = 0.5\)--\(0.6\) is more robust (gives \(O(t^{-0.5})\) to
\(O(t^{-0.6})\) convergence).

\textbf{Proposition 11.2 (Polyak-Ruppert Averaging).} With
Polyak-Ruppert averaging
\(\bar{\theta}_t = \frac{1}{t}\sum_{i=1}^t \theta_i\), the rate improves
to:

\[\mathbb{E}[\|\bar{\theta}_t - \theta^*\|^2] = O(t^{-1}),\]

for \textbf{any} \(a \in (0.5, 1)\), achieving the optimal \(1/t\) rate
without precise tuning (Polyak \& Juditsky, 1992).

\textbf{Status: PROVEN} (standard result, applies to SCX student).

\subsubsection{\texorpdfstring{4.3 Dependence on Feature Dimension
\(d_\phi\)}{4.3 Dependence on Feature Dimension d\_\textbackslash phi}}\label{dependence-on-feature-dimension-d_phi}

The feature dimension \(d_\phi\) affects the convergence rate indirectly
through:

\begin{enumerate}
\def\labelenumi{\arabic{enumi}.}
\item
  \textbf{Covering number of memory bank states} (C1'): The number of
  \(\varepsilon\)-distinguishable memory configurations is
  \(\exp(O((L_S/\varepsilon)^{d_\phi}))\). Larger \(d_\phi\) means more
  possible states, slowing memory bank stabilization.
\item
  \textbf{Gradient variance}: For overparameterized models,
  \(\sigma_\xi^2 = O(d_\theta)\) (gradient noise scales with dimension).
  But for SCX, the effective dimension is \(d_\phi\) (feature dimension
  of the gatekeeper), not \(d_\theta\) (NEP parameter count).
\item
  \textbf{State estimation error}: The clustering error in state
  discovery (Theorem 5) scales as \(O(d_\phi / N_t)\).
\end{enumerate}

\textbf{Proposition 11.3 (Dimension Dependence --- CONJECTURED).} The
convergence rate degrades with \(d_\phi\) as:

\[\mathbb{E}[\|\theta_t - \theta^*\|^2] \leq C \cdot t^{-a} \cdot (1 + \gamma \cdot d_\phi \cdot \log t),\]

where \(\gamma = O(1)\) is a problem-dependent constant. The logarithmic
dependence comes from the covering number of the feature space (metric
entropy \(O(d_\phi \log(1/\varepsilon))\)).

\emph{Status: \textbf{CONJECTURED.} The logarithmic dependence on
\(d_\phi\) is typical for parametric models (Document 07, Proposition
SE-5). The exact interplay with the SGD rate requires further analysis.}

\subsubsection{\texorpdfstring{4.4 Dependence on Noise Rate
\(\eta\)}{4.4 Dependence on Noise Rate \textbackslash eta}}\label{dependence-on-noise-rate-eta}

The noise rate \(\eta\) affects the convergence rate through the
\textbf{signal-to-noise ratio} of the consensus score:

\textbf{Proposition 11.4 (Noise Rate Dependence --- CONJECTURED).} The
effective strong convexity parameter degrades with noise:

\[\mu_{\text{eff}} = \mu \cdot \frac{\Delta_{\min}^2}{\Delta_{\min}^2 + \eta \cdot \sigma_{\text{noise}}^2},\]

where \(\Delta_{\min} = \min_s \Delta_s\) is the state separation gap
(Theorem 1). For high noise (\(\eta \to 1\)) or small separation
(\(\Delta_{\min} \to 0\)), \(\mu_{\text{eff}} \to 0\) and convergence
becomes arbitrarily slow.

\begin{center}\rule{0.5\linewidth}{0.5pt}\end{center}

\subsection{5. Rate Under Polyak-Łojasiewicz
Condition}\label{rate-under-polyak-ux142ojasiewicz-condition}

\subsubsection{5.1 The PL Condition}\label{the-pl-condition}

Many neural network losses satisfy the weaker \textbf{Polyak-Łojasiewicz
(PL)} condition instead of strong convexity:

\[\frac{1}{2} \|\nabla \bar{L}(\theta)\|^2 \geq \mu_{PL} (\bar{L}(\theta) - \bar{L}(\theta^*)), \quad \forall \theta.\]

This is sufficient for linear convergence of gradient descent (Karimi et
al., 2016).

\subsubsection{5.2 PL Rate for SCX
Student}\label{pl-rate-for-scx-student}

\textbf{Theorem 11.4 (Student Rate Under PL --- PROVEN).} Under the PL
condition with constant \(\mu_{PL} > 0\), C2, and C4 with
\(\alpha_t = \alpha_0 t^{-a}\):

\[\boxed{\;\mathbb{E}[\bar{L}(\theta_t) - \bar{L}(\theta^*)] \leq C_{PL} \cdot t^{-a}\;},\]

where \(C_{PL} = \frac{L_g G^2 \alpha_0}{2\mu_{PL}}\).

\emph{Proof.} Standard SGD analysis under PL (Karimi et al., 2016). The
PL condition gives:

\[\mathbb{E}[\bar{L}(\theta_{t+1})] \leq \bar{L}(\theta_t) - \alpha_t \|\nabla \bar{L}(\theta_t)\|^2 + \frac{L_g \alpha_t^2 G^2}{2}.\]

Using \(\|\nabla \bar{L}\|^2 \geq 2\mu_{PL}(\bar{L} - \bar{L}^*)\):

\[\mathbb{E}[\bar{L}(\theta_{t+1}) - \bar{L}^*] \leq (1 - 2\mu_{PL}\alpha_t)(\bar{L}(\theta_t) - \bar{L}^*) + \frac{L_g \alpha_t^2 G^2}{2}.\]

This recurrence gives \(O(t^{-a})\) convergence. \(\square\)

\textbf{Status: PROVEN.} The PL condition is more realistic than strong
convexity for neural networks.

\begin{center}\rule{0.5\linewidth}{0.5pt}\end{center}

\subsection{6. Rate Under General
Non-Convexity}\label{rate-under-general-non-convexity}

\subsubsection{6.1 The Non-Convex Case}\label{the-non-convex-case}

For general non-convex losses (typical deep NEP students), we cannot
guarantee convergence to a global minimum. Instead, we characterize
convergence to a \textbf{stationary point}:

\textbf{Theorem 11.5 (Stationarity Rate, Non-Convex --- PROVEN).} Under
C2, C4 (RM rates), C5, and bounded gradient noise:

\[\boxed{\;\min_{0 \leq i \leq t} \mathbb{E}[\|\nabla \bar{L}(\theta_i)\|^2] \leq \frac{C_{nc}}{t^{1-a}}\;},\]

where \(a \in (0.5, 1)\) is the learning rate exponent, and
\(C_{nc} = \frac{\bar{L}(\theta_0) - \bar{L}_{\inf} + L_g G^2 \sum_{i=0}^\infty \alpha_i^2 / 2}{\sum_{i=0}^t \alpha_i}\).

For \(\alpha_t = \alpha_0 t^{-a}\):

\[\sum_{i=0}^t \alpha_i = \Theta(t^{1-a}),\]

so the rate is \(O(t^{-(1-a)})\). For \(a = 0.5\) (the fastest
RM-compatible rate), this gives \(O(t^{-0.5})\).

\emph{Proof.} Standard SGD stationarity analysis (Ghadimi \& Lan, 2013).
From the \(L_g\)-smoothness:

\[\mathbb{E}[\bar{L}(\theta_{t+1})] \leq \bar{L}(\theta_t) - \alpha_t \|\nabla \bar{L}(\theta_t)\|^2 + \frac{L_g \alpha_t^2 G^2}{2}.\]

Telescoping and rearranging:

\[\sum_{i=0}^t \alpha_i \mathbb{E}[\|\nabla \bar{L}(\theta_i)\|^2] \leq \bar{L}(\theta_0) - \bar{L}_{\inf} + \frac{L_g G^2}{2} \sum_{i=0}^t \alpha_i^2.\]

The minimum over \(i \leq t\) is at most the weighted average, giving
the claimed bound. \(\square\)

\textbf{Status: PROVEN.}

\subsubsection{6.2 Stationarity
vs.~Optimality}\label{stationarity-vs.-optimality}

The rate \(O(t^{-(1-a)})\) guarantees that the gradient norm converges
to zero, \textbf{not} that the parameters converge to \(\theta^*\). In
non-convex landscapes, \(\|\nabla \bar{L}(\theta)\| \to 0\) is
consistent with convergence to a saddle point or a poor local minimum.

\textbf{Proposition 11.5 (Escaping Saddles --- CONJECTURED).} With SGD
noise \(\sigma_\xi^2 > 0\), the probability of converging to a strict
saddle point is zero, and the rate of escape from a saddle with negative
curvature \(\lambda_{\min} < 0\) is exponential in the noise level:

\[\mathbb{P}(\text{escape in } \leq \tau \text{ steps}) \geq 1 - \exp\!\left(-\frac{|\lambda_{\min}| \tau \alpha_t}{\sigma_\xi^2}\right).\]

\emph{Status: \textbf{CONJECTURED} (based on Jin et al., 2017, for
gradient descent with noise).}

\begin{center}\rule{0.5\linewidth}{0.5pt}\end{center}

\subsection{7. Finite-Time (Non-Asymptotic)
Bound}\label{finite-time-non-asymptotic-bound}

\subsubsection{7.1 Why Finite-Time?}\label{why-finite-time}

Asymptotic rates \((t \to \infty)\) are informative about long-run
behavior but say little about what happens at practical timescales
(\(t = 10^2\)--\(10^4\)). A finite-time bound with explicit constants is
needed for practical guarantees.

\subsubsection{7.2 Finite-Time Bound Under Strong
Convexity}\label{finite-time-bound-under-strong-convexity}

\textbf{Theorem 11.6 (Finite-Time Student Bound --- PROVEN).} Under SC
(strong convexity), C2, C4 with \(\alpha_t = \alpha_0 t^{-a}\), for any
\(t \geq 1\) and \(\delta \in (0, 1)\):

\[\boxed{\;\mathbb{P}\!\left(\|\theta_t - \theta^*\|^2 \leq \frac{C_V}{t^a} + \frac{C_{\delta}}{\sqrt{t}} \cdot \sqrt{\log\frac{2}{\delta}}\right) \geq 1 - \delta\;},\]

where:

\[\begin{aligned}
C_V &= \frac{2 G^2 \alpha_0}{\mu} \quad \text{(variance term)}, \\
C_{\delta} &= \frac{4 G \alpha_0}{\mu} \cdot \sqrt{\frac{\sigma_\xi^2}{\mu}} \quad \text{(concentration term)}.
\end{aligned}\]

\emph{Proof.} From the recurrence
\(V_{t+1} \leq (1 - 2\mu\alpha_t) V_t + \alpha_t^2 G^2\), unrolling to
step \(t\):

\[V_t \leq V_0 \prod_{i=0}^{t-1}(1 - 2\mu\alpha_i) + G^2 \sum_{i=0}^{t-1} \alpha_i^2 \prod_{j=i+1}^{t-1} (1 - 2\mu\alpha_j).\]

With \(\alpha_i = \alpha_0 i^{-a}\), the product
\(\prod (1 - 2\mu\alpha_i) \approx \exp(-2\mu \sum \alpha_i) \approx \exp(-c t^{1-a})\).
The dominant term is the variance accumulation, giving \(O(t^{-a})\).
The concentration term uses McDiarmid's inequality on the martingale
difference sequence. \(\square\)

\textbf{Status: PROVEN.} This gives a non-asymptotic guarantee with
explicit dependence on all parameters.

\subsubsection{7.3 Finite-Time Gatekeeper
Bound}\label{finite-time-gatekeeper-bound}

\textbf{Theorem 11.7 (Finite-Time Gatekeeper Bound --- CONJECTURED).}
Under C3, C6', C7, C9, for any \(t \geq 1\) and \(\delta \in (0, 1)\):

\[\boxed{\;\mathbb{P}\!\left(\|S_t - S^*\|_{M_0}^2 \leq \frac{C_S}{t^{b}} + \frac{C_{\delta,S}}{\sqrt{N_t}} \cdot \sqrt{\log\frac{2}{\delta}}\right) \geq 1 - \delta\;},\]

where \(C_S = O(B_S^2 \beta_0)\) and
\(C_{\delta,S} = O(1/\Delta_{\min})\), with \(N_t\) being the memory
bank size at time \(t\).

\emph{Proof sketch.} The gatekeeper's convergence has two components:
(i) the iterative refinement at rate \(t^{-b}\), and (ii) the
statistical error from finite memory, which decays as \(1/\sqrt{N_t}\).
Under monotonic memory growth (B1), \(N_t = \Omega(t)\), so the
statistical error is \(O(1/\sqrt{t})\).

\textbf{Status: CONJECTURED.} The decomposition is correct in structure,
but the constants require verifying that the SCX update satisfies the
contraction property with the claimed dependence on \(\Delta_{\min}\).

\subsubsection{7.4 Combined Finite-Time
Bound}\label{combined-finite-time-bound}

\textbf{Corollary 11.2 (Combined Finite-Time Bound --- CONJECTURED).}
Under the conditions of Theorems 11.6 and 11.7, with \(a = 0.5\) and
\(b = 0.8\), for \(t \geq t_0\):

\[\boxed{\;\mathbb{E}\bigl[\Phi(S_t, \theta_t) - \Phi(S^*, \theta^*)\bigr] \leq \frac{C_\Phi}{\sqrt{t}} \cdot (1 + o(1))\;},\]

where \(C_\Phi\) is the sum of the student and gatekeeper constants.

The \(O(1/\sqrt{t})\) rate is the \textbf{typical stochastic
approximation rate} --- it matches the minimax optimal rate for
stochastic optimization with noisy gradients (Agarwal et al., 2012).

\begin{center}\rule{0.5\linewidth}{0.5pt}\end{center}

\subsection{8. Two-Timescale Rate
Decomposition}\label{two-timescale-rate-decomposition}

\subsubsection{8.1 Fast Timescale: Student
Convergence}\label{fast-timescale-student-convergence}

On the fast timescale, the student sees the gatekeeper as approximately
fixed. The student rate (from Section 3) is:

\[R_\theta(t) = \mathbb{E}[\|\theta_t - \theta^*(S_t)\|^2] = O(t^{-a}).\]

\subsubsection{8.2 Slow Timescale: Gatekeeper
Convergence}\label{slow-timescale-gatekeeper-convergence}

On the slow timescale, the student is approximately at
\(\theta^*(S_t)\), and the gatekeeper evolves according to a reduced
dynamics:

\[S_{t+1} = S_t + \beta_t \cdot \Delta(S_t, \theta^*(S_t)) + \text{noise}.\]

The gatekeeper rate is:

\[R_S(t) = \mathbb{E}[\|S_t - S^*\|^2] = O(t^{-b}) = o(t^{-a}).\]

\subsubsection{8.3 Total Rate}\label{total-rate}

The total convergence rate is the slower of the two:

\[\boxed{\;R_{\text{total}}(t) = \max(R_\theta(t), R_S(t)) = O(t^{-a})\;},\]

since \(b > a\). The \textbf{student bottleneck dominates} --- the
system converges only as fast as the NEP learns.

\subsubsection{8.4 Optimal Rate
Allocation}\label{optimal-rate-allocation}

Given a total ``learning budget'' \(\sum (\alpha_t + \beta_t)\), how
should we allocate between student and gatekeeper?

\textbf{Proposition 11.6 (Optimal Rate Allocation --- CONJECTURED).} The
optimal allocation makes the two rates equal:

\[t^{-a} = t^{-b} \implies a = b.\]

But this \textbf{violates} the two-timescale condition C6'
(\(\beta_t = o(\alpha_t)\), which requires \(b > a\)). Therefore, the
optimal allocation under the two-timescale constraint is:

\[b = a + \varepsilon, \quad \varepsilon \to 0^+,\]

which drives the gatekeeper toward the same rate as the student while
(just barely) maintaining timescale separation.

In practice, setting \(a = 0.5\), \(b = 0.6\) gives near-optimal
convergence while ensuring stable two-timescale behavior.

\begin{center}\rule{0.5\linewidth}{0.5pt}\end{center}

\subsection{9. Comparison with Known
Rates}\label{comparison-with-known-rates}

\subsubsection{9.1 Comparison Table}\label{comparison-table}

\begin{longtable}[]{@{}
  >{\raggedright\arraybackslash}p{(\linewidth - 6\tabcolsep) * \real{0.2222}}
  >{\raggedright\arraybackslash}p{(\linewidth - 6\tabcolsep) * \real{0.3056}}
  >{\raggedright\arraybackslash}p{(\linewidth - 6\tabcolsep) * \real{0.1667}}
  >{\raggedright\arraybackslash}p{(\linewidth - 6\tabcolsep) * \real{0.3056}}@{}}
\toprule\noalign{}
\begin{minipage}[b]{\linewidth}\raggedright
Method
\end{minipage} & \begin{minipage}[b]{\linewidth}\raggedright
Assumption
\end{minipage} & \begin{minipage}[b]{\linewidth}\raggedright
Rate
\end{minipage} & \begin{minipage}[b]{\linewidth}\raggedright
Reference
\end{minipage} \\
\midrule\noalign{}
\endhead
\bottomrule\noalign{}
\endlastfoot
\textbf{SGD (strongly convex)} & \(\mu > 0\) & \(O(t^{-1})\) & Bach \&
Moulines (2011) \\
\textbf{SGD (PL condition)} & PL constant \(\mu_{PL}\) & \(O(t^{-1})\) &
Karimi et al.~(2016) \\
\textbf{SGD (non-convex)} & \(L_g\)-smooth & \(O(t^{-0.5})\)
(stationarity) & Ghadimi \& Lan (2013) \\
\textbf{SGD + averaging} & \(\mu > 0\) & \(O(t^{-1})\) & Polyak \&
Juditsky (1992) \\
\textbf{Two-timescale SA} & Timescale separation & \(O(t^{-a})\) (slower
rate) & Borkar (2008) \\
\textbf{SCX Student (this work)} & \(\mu > 0\) or PL & \(O(t^{-a})\),
\(a \in (0.5, 1]\) & Theorem 11.3 \\
\textbf{SCX Gatekeeper (this work)} & Contraction \(L_S < 1\) &
\(O(t^{-b})\) & Theorem 11.2 \\
\textbf{SCX Coupled (this work)} & Two-timescale & \(O(t^{-a})\) &
Theorem 11.3 \\
\end{longtable}

\subsubsection{9.2 Key Observations}\label{key-observations}

\begin{enumerate}
\def\labelenumi{\arabic{enumi}.}
\item
  \textbf{SCX is not slower than standard SGD}: Under strong convexity,
  the SCX student achieves the same \(O(t^{-a})\) rate as standard SGD.
  The coupling to the gatekeeper does \textbf{not} degrade the
  asymptotic rate (under two-timescale separation).
\item
  \textbf{The gatekeeper converges faster}: Since \(b > a\), the
  gatekeeper converges faster than the student. The system's bottleneck
  is the NEP training, not the SCX calibration.
\item
  \textbf{Polyak averaging provides \(O(1/t)\)}: Polyak-Ruppert
  averaging of the student iterates achieves the optimal \(O(1/t)\) rate
  without degrading the gatekeeper convergence.
\end{enumerate}

\begin{center}\rule{0.5\linewidth}{0.5pt}\end{center}

\subsection{10. Gap Between Upper and Lower
Bounds}\label{gap-between-upper-and-lower-bounds}

\subsubsection{10.1 Information-Theoretic Lower
Bound}\label{information-theoretic-lower-bound}

\textbf{Theorem 11.8 (Minimax Lower Bound for Student Convergence ---
PROVEN).} For any algorithm that updates \(\theta_t\) using stochastic
gradients with variance \(\sigma_\xi^2 > 0\), under \(\mu\)-strong
convexity:

\[\mathbb{E}[\|\theta_t - \theta^*\|^2] \geq \frac{\sigma_\xi^2}{\mu^2} \cdot \frac{1}{t}.\]

\emph{Proof.} Standard minimax lower bound for stochastic optimization
(Agarwal et al., 2012; Raginsky \& Rakhlin, 2011). The Cramér-Rao bound
for sequential estimation gives \(\Omega(1/t)\). \(\square\)

\subsubsection{10.2 Gap Analysis}\label{gap-analysis}

\begin{longtable}[]{@{}
  >{\raggedright\arraybackslash}p{(\linewidth - 6\tabcolsep) * \real{0.1778}}
  >{\raggedright\arraybackslash}p{(\linewidth - 6\tabcolsep) * \real{0.2667}}
  >{\raggedright\arraybackslash}p{(\linewidth - 6\tabcolsep) * \real{0.2889}}
  >{\raggedright\arraybackslash}p{(\linewidth - 6\tabcolsep) * \real{0.2667}}@{}}
\toprule\noalign{}
\begin{minipage}[b]{\linewidth}\raggedright
Regime
\end{minipage} & \begin{minipage}[b]{\linewidth}\raggedright
Upper Bound
\end{minipage} & \begin{minipage}[b]{\linewidth}\raggedright
Lower Bound
\end{minipage} & \begin{minipage}[b]{\linewidth}\raggedright
Gap Factor
\end{minipage} \\
\midrule\noalign{}
\endhead
\bottomrule\noalign{}
\endlastfoot
\(\alpha_t = t^{-1}\), SC & \(O(t^{-1})\) & \(\Omega(t^{-1})\) &
\(O(1)\) --- \textbf{tight} \\
\(\alpha_t = t^{-0.5}\), SC & \(O(t^{-0.5})\) & \(\Omega(t^{-1})\) &
\(O(t^{0.5})\) --- \textbf{loose} \\
Non-convex, stationarity & \(O(t^{-0.5})\) & N/A (no optimality
guarantee) & --- \\
Gatekeeper, contraction & \(O(t^{-b})\) &
\(\Omega(1/N_t) = \Omega(t^{-1})\) & \(O(t^{1-b})\) --- \textbf{loose
for \(b < 1\)} \\
\end{longtable}

\textbf{Interpretation}: With optimal tuning (\(a = 1\), Polyak
averaging), the student rate matches the minimax lower bound up to
constants. With suboptimal tuning (\(a = 0.5\)), the rate is suboptimal
by \(O(\sqrt{t})\). The gatekeeper rate has room for improvement --- the
\(t^{-b}\) dependence from iterative refinement is looser than the
\(1/\sqrt{N_t}\) statistical rate from memory accumulation.

\subsubsection{10.3 Tightening the Gatekeeper
Rate}\label{tightening-the-gatekeeper-rate}

\textbf{Proposition 11.7 (Improved Gatekeeper Rate via Batch Updates ---
CONJECTURED).} If the gatekeeper is updated using \textbf{full-batch}
SCXUpdate on the entire memory bank (rather than incremental updates),
its convergence rate improves to:

\[\mathbb{E}[\|S_t - S^*\|_{M_0}^2] \leq \frac{C_S'}{\sqrt{N_t}},\]

matching the statistical rate. Since \(N_t = \Theta(t)\) (constant batch
sizes), this gives \(O(t^{-0.5})\).

\emph{Status: \textbf{CONJECTURED.} The full-batch update eliminates the
iterative refinement bottleneck, leaving only the statistical error from
finite memory. This is faster than the incremental rate when
\(b < 0.5\).}

\begin{center}\rule{0.5\linewidth}{0.5pt}\end{center}

\subsection{11. Summary}\label{summary}

\subsubsection{11.1 Proven Results}\label{proven-results}

\begin{longtable}[]{@{}
  >{\raggedright\arraybackslash}p{(\linewidth - 4\tabcolsep) * \real{0.3636}}
  >{\raggedright\arraybackslash}p{(\linewidth - 4\tabcolsep) * \real{0.3636}}
  >{\raggedright\arraybackslash}p{(\linewidth - 4\tabcolsep) * \real{0.2727}}@{}}
\toprule\noalign{}
\begin{minipage}[b]{\linewidth}\raggedright
Result
\end{minipage} & \begin{minipage}[b]{\linewidth}\raggedright
Status
\end{minipage} & \begin{minipage}[b]{\linewidth}\raggedright
Rate
\end{minipage} \\
\midrule\noalign{}
\endhead
\bottomrule\noalign{}
\endlastfoot
Student rate under strong convexity (stationary dist.) & \textbf{PROVEN}
& \(O(t^{-a})\), \(a \in (0.5, 1]\) \\
Student rate under PL condition & \textbf{PROVEN} & \(O(t^{-a})\) \\
Student stationarity rate (non-convex) & \textbf{PROVEN} &
\(O(t^{-(1-a)})\) \\
Finite-time student bound & \textbf{PROVEN} &
\(O(t^{-a}) + O(t^{-0.5})\) (concentration) \\
Minimax lower bound & \textbf{PROVEN} & \(\Omega(t^{-1})\) \\
\end{longtable}

\subsubsection{11.2 Conjectured Results}\label{conjectured-results}

\begin{longtable}[]{@{}
  >{\raggedright\arraybackslash}p{(\linewidth - 4\tabcolsep) * \real{0.2581}}
  >{\raggedright\arraybackslash}p{(\linewidth - 4\tabcolsep) * \real{0.2581}}
  >{\raggedright\arraybackslash}p{(\linewidth - 4\tabcolsep) * \real{0.4839}}@{}}
\toprule\noalign{}
\begin{minipage}[b]{\linewidth}\raggedright
Result
\end{minipage} & \begin{minipage}[b]{\linewidth}\raggedright
Status
\end{minipage} & \begin{minipage}[b]{\linewidth}\raggedright
Expected Rate
\end{minipage} \\
\midrule\noalign{}
\endhead
\bottomrule\noalign{}
\endlastfoot
Coupled rate under strong convexity & \textbf{CONJECTURED} &
\(O(t^{-a})\) \\
Gatekeeper rate & \textbf{CONJECTURED} & \(O(t^{-b})\) \\
Finite-time gatekeeper bound & \textbf{CONJECTURED} &
\(O(t^{-b}) + O(t^{-0.5})\) \\
Optimal rate allocation \(b = a + \varepsilon\) & \textbf{CONJECTURED} &
\(O(t^{-a})\) \\
\end{longtable}

\subsubsection{11.3 Parameter Dependence
Summary}\label{parameter-dependence-summary}

\begin{longtable}[]{@{}
  >{\raggedright\arraybackslash}p{(\linewidth - 4\tabcolsep) * \real{0.1864}}
  >{\raggedright\arraybackslash}p{(\linewidth - 4\tabcolsep) * \real{0.4068}}
  >{\raggedright\arraybackslash}p{(\linewidth - 4\tabcolsep) * \real{0.4068}}@{}}
\toprule\noalign{}
\begin{minipage}[b]{\linewidth}\raggedright
Parameter
\end{minipage} & \begin{minipage}[b]{\linewidth}\raggedright
Effect on Rate Constant
\end{minipage} & \begin{minipage}[b]{\linewidth}\raggedright
Effect on Rate Exponent
\end{minipage} \\
\midrule\noalign{}
\endhead
\bottomrule\noalign{}
\endlastfoot
\(L_f\) (Lipschitz) & \(C \propto L_f^3\) & None \\
\(\alpha_t = t^{-a}\) & \(C \propto \alpha_0\) & Rate = \(t^{-a}\) \\
\(d_\phi\) (dimension) & \(C \propto 1 + \gamma d_\phi \log t\) & None
(logarithmic) \\
\(\eta\) (noise rate) & \(C \propto 1/\Delta_{\min}^2(\eta)\) & None \\
\(\mu\) (curvature) & \(C \propto 1/\mu\) & None \\
\(\sigma_\xi^2\) (gradient noise) & \(C \propto \sigma_\xi^2\) & None \\
\end{longtable}

\subsubsection{11.4 Practical
Recommendations}\label{practical-recommendations}

\begin{enumerate}
\def\labelenumi{\arabic{enumi}.}
\tightlist
\item
  \textbf{Use Polyak-Ruppert averaging} on the student iterates to
  achieve \(O(1/t)\) rate regardless of \(a \in (0.5, 1)\).
\item
  \textbf{Set \(a = 0.5\), \(b = 0.6\)} for robust convergence with
  near-optimal rate allocation.
\item
  \textbf{Use batch gatekeeper updates} if computationally feasible, to
  achieve the \(O(1/\sqrt{N_t})\) statistical rate.
\item
  \textbf{Monitor \(\|\nabla \bar{L}(\theta_t)\|\)} as a practical
  convergence diagnostic --- it should decay as \(t^{-(1-a)/2}\).
\item
  \textbf{The feature dimension \(d_\phi\) has only logarithmic effect}
  on the asymptotic rate --- scaling to high-dimensional features is
  feasible.
\end{enumerate}

\begin{center}\rule{0.5\linewidth}{0.5pt}\end{center}

\emph{End of 11\_convergence\_rate.md --- Convergence Rate Analysis}

\begin{center}\rule{0.5\linewidth}{0.5pt}\end{center}

\subsection{References}\label{references}

\begin{enumerate}
\def\labelenumi{\arabic{enumi}.}
\tightlist
\item
  Bach, F., \& Moulines, E. (2011). Non-asymptotic analysis of
  stochastic approximation algorithms for machine learning. \emph{NIPS
  2011}.
\item
  Polyak, B. T., \& Juditsky, A. B. (1992). Acceleration of stochastic
  approximation by averaging. \emph{SIAM Journal on Control and
  Optimization}, 30(4), 838-855.
\item
  Ghadimi, S., \& Lan, G. (2013). Stochastic first- and zeroth-order
  methods for nonconvex stochastic programming. \emph{SIAM Journal on
  Optimization}, 23(4), 2341-2368.
\item
  Karimi, H., Nutini, J., \& Schmidt, M. (2016). Linear convergence of
  gradient and proximal-gradient methods under the Polyak-Łojasiewicz
  condition. \emph{ECML PKDD 2016}.
\item
  Needell, D., Srebro, N., \& Ward, R. (2014). Stochastic gradient
  descent, weighted sampling, and the randomized Kaczmarz algorithm.
  \emph{NIPS 2014}.
\item
  Borkar, V. S. (2008). \emph{Stochastic Approximation: A Dynamical
  Systems Viewpoint}. Cambridge University Press.
\item
  Agarwal, A., et al.~(2012). Information-theoretic lower bounds on the
  oracle complexity of stochastic convex optimization. \emph{IEEE
  Transactions on Information Theory}, 58(5), 3235-3249.
\item
  Raginsky, M., \& Rakhlin, A. (2011). Information-based complexity,
  feedback and dynamics in convex programming. \emph{IEEE Transactions
  on Information Theory}, 57(10), 7036-7056.
\item
  Jin, C., et al.~(2017). How to escape saddle points efficiently.
  \emph{ICML 2017}.
\item
  Kushner, H. J., \& Yin, G. G. (2003). \emph{Stochastic Approximation
  and Recursive Algorithms and Applications} (2nd ed.). Springer.
\end{enumerate}

\end{document}
